\documentclass[a4paper]{article}

\usepackage{amsmath}
\usepackage{geometry}
\usepackage{amsfonts}
\usepackage{amssymb}
\usepackage{graphicx}

% \usepackage{parskip}
\usepackage[magyar]{babel}
\usepackage[T1]{fontenc}

\geometry{hmargin = 3cm, top = 3cm, bottom = 2.5cm, headsep = 0.9cm}

\begin{document}
\pagenumbering{gobble}

\begin{center}
\Large{Szomszédvezérelt integrálható spin-1 láncok klasszifikációja}\\
\hspace{10pt}

Gulyás Kristóf $^\dagger$, fizika BSc 3. év \\
\normalsize
Témavezető: Pozsgay Balázs $^\ast$ \\
\normalsize
% \large

\hspace{10pt}

\small{$^\dagger$ kristof0905@gmail.com\\}
\small{$^\ast$ pozsgaybalazs@gmail.com\\}
\hspace{5pt}
\end{center}


{\large
\begin{abstract}
A kvantumosan \textit{integrálható} modellek -- melyek végtelen sok lokális megmaradó mennyiséggel rendelkeznek, és sok esetben egzaktul megoldhatók -- kiemelt szerepet játszanak többek közt a zárt kvantumrendszerek nemegyensúlyi dinamikájának megértésében, numerikus számolások validálásában és kvantumhardverek benchmarkolásában. Az általános konstrukcióban a Yang--Baxter-egyenlet biztosítja a szükséges algebrai struktúrát: az egyenlet egy megoldásából -- az úgynevezett $R$-mátrixból -- eltolásinvariáns láncok Hamilton-sűrűsége származtatható. Fordított irányban egy adott Hamilton-sűrűség integrálhatósága a \textit{Reshetikhin-feltétellel}, azaz egy speciális megmaradó mennyiség létezésével tesztelhető.

Munkám során háromdimenziós (spin-1) lokális Hilbert-terekkel rendelkező, háromhelyes, szomszédvezérelt Hamilton-sűrűségű láncokat vizsgáltam. Ez a szerkezet a kétszeresen vezérelt kvantumkapuk logikáját követi, így az ilyen modellek természetes kapcsolatban állnak a kvantumáramkörökkel és a kvantumos sejtautomatákkal. A vizsgált modellek a közepes hatótávolságú integrálható rendszerek családjába tartoznak, melyek tárgyalására jól kidolgozott keretrendszer létezik, azonban ezt az irodalomban főként qubites rendszerekre alkalmazták. Fizikailag motivált szimmetriák betartását követeltem meg a lánctól, majd a már említett Reshetikhin-feltételt egy homogén, harmadfokú polinomiális egyenletrendszerre vezettem vissza. Ezt egy saját fejlesztésű szoftverrel, Gröbner-bázisokon alapuló egzakt algebrai módszerekkel oldottam meg, a lokális bázistranszformációkkal ekvivalens megoldások kiszűrésével. Az algoritmus rendeltetésszerű működését több tétel bizonyításával is alátámasztottam. Klasszifikációm során egyes szimmetriaosztályokban nemtriviális integrálható modelleket azonosítottam, míg másokban kizártam a létezésüket. Az algebrai megoldást nyújtó algoritmust az irodalomban fellelhető modellek (például a spin-1/2 szomszédvezérelt láncok) reprodukálásával is validáltam.
    
\end{abstract}
}




\end{document}